\documentclass{article}
\usepackage{graphicx, hyperref} % Required for inserting images
\usepackage[dvipsnames]{xcolor}
\usepackage{amssymb}
\usepackage{amsmath}
\usepackage{amsthm}
\usepackage{amsfonts}
\usepackage{bbm}
\newtheorem{theorem}{Theorem}
\newtheorem{observation}[theorem]{Observation}
\newtheorem{lemma}{Lemma}
\newtheorem{definition}{Definition}

\usepackage{tikz}
\usetikzlibrary{bayesnet}
\usepackage{url}
\usepackage{booktabs}
\tikzset{%
  zeroarrow/.style = {-stealth,dashed},
  onearrow/.style = {-stealth,solid},
  c/.style = {circle,draw,solid,minimum width=2em,
        minimum height=2em},
  r/.style = {rectangle,draw,solid,minimum width=2em,
        minimum height=2em}
}

\newcommand{\midlinewidth}{1.0pt}
\newcommand{\incmidlinewidth}{1.7pt}
\newcommand{\middist}{20pt}
\newcommand{\halfdist}{13pt}
\definecolor{petroil2} {RGB} {36, 165, 175}
\definecolor{gold2} {RGB} {255, 130, 0}

\title{Causal Explanations in Chirho}
\author{}
\date{\today}

\newcommand{\poorva}[1]{{\color{blue}[Poorva: #1]}}


\newtheorem{corollary}[theorem]{Corollary}


\begin{document}

\maketitle

% \section{Preliminary Definitions}

% % \poorva{The following definition is very operational, is it posisble to have something declarative}
% % \begin{definition}[Sample Operator]
% % A sample operator \texttt{sample(.)} takes as input a probability distribution and outputs a sample from the specified distribution.
% % \end{definition}

% \begin{definition}[Sample Statement]
% A sample statement \(X := D\) assigns the probability distribution \(D\) to the variable \(X\).
% \end{definition}

% \begin{definition}[Causal Probabilsitic Program]
% A causal probabilistic program \(C\) can be defined as a tuple \((X, F\) where X is a set of variables and F is a set of statements consisting of variables X. These statements can be either sample statements assigning distributions to \(X\) or functional expressions:

% \[X_j := D_j(\mathit{Pa}(X_j)) \;\;\; \mathit{ or } \;\;\; X_j := f(\mathit{Pa}(X_j)) \;\;\; j = 1, 2, \ldots, |X|\]

% where \(\mathit{Pa}(X_j) \subseteq \{X \backslash X_j\}\) are called the parents of \(X\). Note that in causal probabilistic programs, we assume that the parents of a node in \(X\) are its causal parents.
% \end{definition}

% \poorva{To completely comprehend the following definition, we should define what does it mean for some assignment to be the cause of an event E, we can appeal to actual causality I guess. Also define necessary and sufficient (probably from chapter 9 in Pearl's book)}
% \begin{definition}[Causal Explanations]
% Given a causal probabilistic program \(C = (X, F)\), the causal explanation of the event \(E := \tilde{X} = x \;\;\; \tilde{X} \in X\) is defined a minimal set of variables in \(X\) and the values they take that are necessary and sufficient to cause event E. 
% \end{definition}

% \poorva{I understand the goal is to formalize the transformation from causal program to probabilistic program and show that some probabilities match, I believe it would be imperative to define the syntax of the program in that case.}

% \section{Transformation from Search to Inference}

% \poorva{This is a description of the current implementation and need to be refined into a formal description}

% \subsection{SearchForExplanation}

% \texttt{SearchForExplanation} currently takes the following inputs:
% \begin{itemize}
%     \item A causal probabilistic model
%     \item An outcome: \texttt{consequent}
%     \item A list of candidates to be included in the set of causes: \texttt{antecedents}
%     \item Alternate values for \texttt{antecedents}: \texttt{alternatives}
%     \item Other arguments which we would describe later 
% \end{itemize}

% \noindent and outputs a probabilistic model conditioned on the following:
% \begin{itemize}
%     \item \textbf{Sufficiency condition}: intervening on antecedents with their given values result in the outcome
%     \item \textbf{Necessity condition}: intervening on antecedents with their alternative values result in outcome not happening
% \end{itemize}

% \subsection{Example: Forest fire}


% \poorva{Definitions in sequence:
% 1. Structural Causal Models and the probability distribution defined by them
% 2. Sufficiency condition
% 3. Necessity condition
% 4. Factivity condition
% 5. conditioning on above
% 6. Transformation with auxiliary variables (conditional conditioning: side note for Zenna)

% Theorem 1: Prove that existing definitions comply
% Theorem 2: New definitions exist}

\section{Introduction}

\texttt{Chirho} is a causal probabilistic programming language that enables users to write their causal probabilistic models as programs and perform computations of causal queries. Interventions are performed as program transformations agnostic to the inference algorithms and higher-level causal queries are expressed in terms of interventions.

In this paper our focus is on answering queries related to causal explanations, that is, given a causal model, what factors explain the outcome. 


\section{Probability of causation: intuitions}


To motivate the technical decisions and definitions that will follow, let us first think through a simplified example involving responsibility attribution.

Bob applied for a home mortgage loan with the local branch of Old Boys’ Club Bank (OBCB), and his application was denied. The OBCB policy requires applicants to pass a credit score check, but it only performs this check for male applicants; all non-male applicants are rejected outright. Assuming this policy holds, why was Bob’s loan application denied? An intuitive explanation is that his credit score was insufficient—had his score been higher, the outcome would have been different. In contrast, his gender does not appear to be responsible: had it been different, his application still would have been denied.


% \tikz{ %
%         \node[latent] (g) {Gender} ; %
%         \node[latent, right=of g] (cs) {Credit Score} ;
%         \node[latent, left=of cs, right=of g, below=of g] (ccf) {Credit Check Failed} ;
%         \node[latent] (aa, yshift=0.5cm) {Application Accepted}
%         % \node[latent, left=of y, yshift=0.5cm] (mu) {$\mu$} ; %
%         % \node[latent, left=of y, yshift=-0.5cm] (tau) {$\tau$} ; %
%         \edge {} {} ; %
%       }

This suggests a counterfactual approach to feature responsibility attribution: we can identify the features responsible for a particular outcome by asking counterfactual questions—would the outcome have been different if a specific feature had a different value? This is called the \textit{but-for clause}: Bob wouldn’t have been denied the loan but for his insufficient credit score.  

To formalize this, suppose we are dealing with a model $M$ in a particular run $u$, where $A=a$ and $B=b$ hold, that is, $(M, u) \models A=a, B=b$.  
A first attempt at defining the relevant counterfactual notion of responsibility might be:
\begin{align*}
A=a \text{ is causally responsible for } B=b \text{ in } (M, u) \text{ just in case } \\   (M, u)[A \leftarrow a'] \not \models B=b
\end{align*}
for some $a' \neq a$. The arrow represents an intervention on $A$, setting it to an alternative value. Bob’s credit score satisfies this condition: had his credit score been higher, the loan might have been approved.  

However, this notion turns out to be insufficiently fine-grained. Consider a slightly more complicated case: Alice, a woman with a credit score that would have been insufficient if checked, applies with OBCB for the same loan. Her application was rejected, but for a different reason---her gender alone was sufficient for rejection, and no credit score check was performed. Why was Alice denied the loan?  

Under the but-for clause, neither her gender nor her credit score individually appears responsible. No matter how high her credit score, her application would still have been denied because the bank does not evaluate credit scores for non-male applicants. Similarly, her gender alone does not fully explain the outcome—male applicants are denied based on their credit score if it is too low, so had her gender been different, her application would still have been rejected (as her credit score was insufficient anyway).

One way to address such cases is to generalize to allow sets or vectors of variables to be jointly responsible:
\begin{align*}
\vec{X} = \vec{x} \text{ is causally responsible for } B=b \text{ in } (M, u) \text{ just in case } \\  (M, u)[\vec{X} \leftarrow \vec{x'}] \not \models B=b
\end{align*}
for some $\vec{x'}$ everywhere different from $\vec{x}$, where no proper subset of $\vec{X}$ has this property (to avoid attributing responsibility to irrelevant properties).  

Under this generalization, the set $(\text{gender = female, credit score = any value})$ is jointly responsible for Alice's application being denied. Changing these variables together---e.g., setting $(\text{gender = male, credit score = sufficient})$—would change the decision. Thus, her denial was caused by the interplay of both features.

But what can we say about the individual features? Neither Alice’s gender nor her credit score is individually deemed responsible. One possible reaction is to say that the whole set bears responsibility, even though the features taken separately do not. Another is to argue that any feature participating in a responsible set is partially accountable by the same token.  

Both approaches have issues, as they imply symmetry: both features are deemed equally responsible. This conflicts with the intuition that since OCBC did not even evaluate Alice’s credit score, her denial was actually caused by her gender, not by her credit score.\footnote{Examples of this sort, involving both over-determination and undercutting of causal impact, abound. Suppose Jim is both stabbed and poisoned and dies as a result. Because stabbing works quickly, it kills Jim before the poison takes effect. In this case, it’s the stabbing that caused his death. Yet analogous but-for clauses fail to break the symmetry here as well.}


As suggested by \cite{actualCausalityHalpern}, a way to break this symmetry is to consider the context. Notably, Alice did not even fail the credit score check. We can imagine extending our model by adding a variable to represent this fact. If we keep this assumption at its factual value, counterfactually, Alice would have been granted the lease if she were male. Asymmetrically, there is no fixed context such that, had her credit score been higher, she would have been granted the lease—her gender blocked the credit check in the first place. This motivates the following formalization:  
\begin{align*}
\vec{X} = \vec{x} \text{ is causally responsible for } B=b \text{ in } (M, u) \text{ just in case }  \\
(M, u)[\vec{X} \leftarrow \vec{x'}, \vec{W} \leftarrow \vec{w^\star}] \not \models B=b,
\end{align*}
\noindent for some \(\vec{x'}\) distinct from \(\vec{x}\), and for some set of context variables \(\vec{W}\) (referred to as "witnesses") with their actual values \(\vec{w^\star}\), where no proper subset of \(\vec{X}\) satisfies this property (this ensures responsibility is not attributed to irrelevant variables).  

This formulation explicates actual causality by evaluating but-for clauses with certain contextual assumptions held fixed. The actual causality claim holds if a context exists where the corresponding restricted but-for clause holds.   To answer causal responsibility queries using this or similar notions, we need to perform two distinct searches: one to identify potential causally responsible sets of variables, and another to determine the appropriate context sets to hold fixed. These searches are key components of our approach.



While we do agree that bringing witness nodes into the picture to break symmetries is a great idea, we will depart from this notion in some respects. For one thing, this definition is formulated in terms of an existential quantifier: it is enough that one actual context being kept fixed validates the counterfactual dependency. We will prefer a more subtle approach where information from counterfactual dependencies across different witness sets is used (and potentially marginalized). After all, the mere existence of a context in which a counterfactual dependency arises might come about in importantly different ways. The context might have been a very rare combination of background factors leading to a very unlikely accidental dependency, or it might have been the presence of obvious factors everyone should expect to be present. This should make a difference to the responsibility attribution.



Moreover, we would like to be probabilistic about responsibility attribution. Consider a modified scenario, in which OBCB proceeds stochastically. They in fact proceed with credit score checks for randomly chosen 20\% women and they randomly drop 10\% male candidates prior to credit check. If a check fails, they nevertheless take the risk for randomly chosen 5\% males, and if it succeeds, they block the lease for randomly chosen 10\% females anyway.

Let's set choosing witnesses and considering feature sets aside, and focus on the counterfactual impact of changing individual features. First, Ben. With his current features, Ben's chance of getting the lease is $.9 \times 0.05 = 0.045$. If he were not male, this chance would decrease to $.2 \times 0=0$.
If he had a sufficient credit score, his probability of getting the lease would jump to .9.  So changing his gender does not increase the probability of the outcome being different (quite the opposite), while changing his credit score would significantly do so. Probabilistic responsibility attribution should be sensitive to such differences.

The probability we discussed above is the \emph{probability of necessity} \cite{pearl2022probabilities}, the probability that the outcome would not have occurred in the absence of the antecedent, given that both occurred: 
$$ PN(X=x, Y=y) = P(Y_{x'} \neq y \vert X=x, Y=y)$$ where $Y_{x'}$ is the value of $Y$ under the intervention setting $X$ to an alternative value $x'\neq x$ of $X$.  In our approach, we will combine the intuition that such probabilities matter in responsibility attribution, except we will do so in a context-aware fashion, making sure search over witness nodes to be preempted is factored in.


What about Alice? Currently, her probability of getting the lease is 0. If she were male, her probability would rise to 0.045, the same as Ben's. If she were to have a sufficient credit score while remaining female, the probability of a different outcome would increase to 0.18. Based solely on the probability of necessity, one might conclude that Alice's gender is barely responsible for the decision and that her credit score is several times less influential than Ben's. This conclusion feels misguided.

In fact, we should think in terms of at least two scenarios. On one hand, we should follow the intuitions that underlie both the but-for clause and the probability of necessity: investigate whether changing the feature would change the outcome. We will call possible worlds in which such scenarios happen \emph{necessity worlds}. On the other hand, we should ask how \emph{sufficient} the given feature is for a given outcome: how often an intervention that mimics the actual state of the candidate would cause the outcome?\footnote{Example:  the probability of necessity of breathing for doing a PhD in computer science is very high, but the corresponding probability of sufficiency is rather low.}   We will call possible worlds in which such interventions happen \emph{sufficiency worlds}. A relevant notion present in the literature is that of the probability of sufficiency, which mirrors that of necessity:
$$ PS(X=x, Y=y) = P(Y_{x} \neq y \vert X=x', Y=y')$$
Coming back to Alice, suppose she isn't female, and her lease was approved. What's the probability of her application being rejected if she were female?
In this scenario, her credit score is still insufficient, so she was just one of the lucky 5\% of male applicants whose credit score has been evaluated. If she were female, the probability of her getting the lease would therefore be 0. So the probability of sufficiency for her gender is 1. This captures the intuition that her gender played a strong role. Alternatively, suppose her credit score is sufficient, she is still female, and she did get the lease. This means she's among the lucky 90\% of the lucky 20\% women whose credit score has been evaluated. The probability of her application being rejected if her credit score was insufficient is 1 --- which, again, aligns with our intuitions that an insufficient credit score is, in the case of a female applicant, completely sufficient for the application failure.\footnote{How about the probabilities of sufficiency for Ben? If he were female and got the lease, the only way this could have happened is by having a sufficient credit score (and the probability of this outcome is 0.18, i.e. the probability of rejection is  .82). Then, his being male makes the probability of rejection drop to .1, which suggests the probability of sufficiency for him being male is low, as expected. Now suppose his credit score is sufficient, and he got the lease. What's the probability of his application being rejected if we intervene to make his credit score insufficient?  That's 1- (.9 * .05) = .955. So the probability of sufficiency for credit score is fairly high.}

The main observation here is that a properly nuanced picture of the causal role of a feature in the context of a binary outcome should be two-sided, capturing how changing the feature could change the outcome, but also, gauging the extent to which having the feature as it is suffices for the outcome. Therefore, in our approach, the joint probability of necessity and sufficiency will be intertwined with context sensitivity and search for causal sets.  In what follows we move to what we consider a natural way of going about this.















\section{Probability of causation: definitions}



The above description  is quite informal and we make it more concrete it within the framework of  Structural Causal Models:

\begin{definition}[Causal Model]
    A causal model is s triple \(M = \langle U, V, F \rangle\) where:
\begin{itemize}
    \item U is a set of background (i.e. exogenous) variables  that are determined by factors outside the model,
    \item V is a set \(\{V_1, V_2, \ldots V_n\}\) of  variables whose values are determined by other variables already in the model (these are called endogenous variables),
    \item F is a set of functions \(\{f_1, f_2, \ldots, f_n\}\) such that each \(f_i\) is a mapping from \(U_i \cup \mathit{Pa}_i\) to \(V_i\) where \(U_i \subseteq U\), \(\mathit{Pa}_i \subseteq V \backslash V_i\) and the entire set forms a mapping from U to V. In other words, for each $v_i$, its associated function  \(f_i\) 
    assigns a value to \(V_i\) that depends on an appropriate set of "upstream" variables in \(V \cup U\): \[v_i = f_i(\mathit{pa}_i, u_i), \;\;\; i = 1, 2, \ldots n\]
    so that the entire set \(F\) has a unique solution \(V(u)\).
\end{itemize}
\end{definition}

Such models can be  equipped with probabilities to define probabilistic causal models as follows~\cite{causalityPearl}:
\begin{definition}[Probabilistic Causal Model]
    A probabilistic causal model is a pair \(\langle M, P(u) \rangle\), where \(M\) is a causal model and \(P(u)\) is a probability function defined over the domain of \(U\).
\end{definition}

This definition allows us to associate probability distributions with causal models. That is, given a distribution $P(u)$ over the exogenous nodes, we have a pushforward joint distribution \(P_j\) on all the model variables \(U \cup V\) as follows:
\begin{align} \label{eq:joint} P_j(\vec{u}, \vec{v}) = P(\vec{u}) \cdot \mathbbm{1}\left ( (M, \vec{u}) \vDash \vec{v}\right)\end{align}
where \(\vec{u}\) and \(\vec{v}\) are assignments to \(U\) and \(V\) respectively. We can also recover marginal probabilities for any model variables $\vec{a} \subseteq U\cup V$ as follows:
\[
P(\vec{a}) = \sum_{\vec{b}} P_j(\vec{a}, \vec{b}), \mbox{  where } \vec{b} = U\cup V \setminus \vec{a}.
\]


\begin{definition}[Causal set]
    Given a probabilistic causal model, \(\langle U, V, F, P(u)\rangle\) a causal set is a set of variable value pairs \(\{\vec{X} = \vec{x} | \vec{X} \subseteq  V, \vec{x} \in \mathit{dom}(\vec{X})\}\). 
\end{definition}

As such, a causal set is just a set of variables and some values; what makes it causal is that such sets will be considered potential causes or will be intervened on. 
We define interventions in causal models in alignment with Pearl's do calculus~\cite{causalityPearl}  as follows:

\begin{definition}[Interventions]
    Given a causal model \(\langle U, V, F \rangle\), variable \(X \in V\) can be intervened with value \(x \in \textit{dom}(X)\) by replacing the function for \(X\) by \(X = x\). We use the notation \([X \leftarrow x]\) to denote the intervention.
\end{definition}

In what follows, causal sets will take on multiple roles. At times, they will represent sets of \emph{potential} causes---whose subsets may be intervened upon---in which case we write them as $\vec{X}$. In other cases, they will serve as \emph{potential} contexts to be kept fixed, denoted by $\vec{W}$. When referring to specific sets we \emph{in fact} intervene on, we use $\vec{C}$; and when referring to context variables that \emph{in fact} are held fixed, we use $\vec{T}$.

\subsection{Conditioning on Sufficiency and Necessity}


The definitions of sufficiency and necessity below are similar to Pearl's, except for two main differences. We do not consider the values of the concerned variables in the factual world. This is mainly because we intend to use these definitions in probabilistic causal models across different possible factual worlds, while it is easy to test factivity separately if need be. We also add context sets $W$ similar to contingencies in Halpern's definition of actual causality~\cite{actualCausalityHalpern} to incorporate the intuition that token-level causality should be context-sensitive.




\begin{definition}[Sufficiency]
    Given a causal model \(M\) and an assignment to exogenous variables \(u\), the causal set \(\{\vec{C} = \vec{c}\}\) with the context set T is \textit{sufficient} for outcome \(\varphi\) if the following holds:
    \begin{align}
    \tag{$S_{\vec{C} = \vec{c}, \vec{T}}(\vec{u}, \varphi)$}
    (M, \vec{u}) 
    \vDash [\vec{C} \leftarrow \vec{c}, \vec{T} \leftarrow \vec{t}^\star]\varphi
    \end{align}
     where $\vec{t}^\star$ are the actual values of $\vec{T}$, that is, \((M, \vec{u}) \vDash \{\vec{T} = \vec{t}^\star\}\). Assuming the model $M$ is held fixed,  we refer to this condition as \(S_{\vec{C} = \vec{c}, \vec{T}}(\vec{u}, \varphi)\).

\end{definition}

\begin{definition}[Necessity]
    Given a causal model \(M\) and an assignment to exogenous variables \(u\), the causal set \(\{\vec{C} = \vec{c}\}\) with the context set T is \textit{necessary} for outcome \(\varphi\) if the following holds for some  \(\vec{c'} \in \mathit{dom}(\vec{C})\setminus \vec{c}\):
    \begin{align} \label{eq:nec}
        \tag{\(N_{\vec{C} = \vec{c}, \vec{T}}(\vec{u}, \varphi)\)}
    (M, \vec{u}) \vDash [\vec{C} \leftarrow \vec{c'}, \vec{T} \leftarrow \vec{t}^\star]\neg \varphi\end{align}
    where \((M, \vec{u}) \vDash \{\vec{T} = \vec{t}^\star\}\).
    %, \vec{x'} \neq \vec{x}\).
    Holding the model fixed, we refer to the necessity condition as \(N_{\vec{C} = \vec{t}, \vec{T}}(\vec{u}, \varphi)\).
\end{definition}

\noindent Note that in the definition of necessity, one can alternatively consider a particular \(\vec{c'}\in \mathit{dom}(C)\) when intending to work with a more explicitly contrastive notion. 

Now, given the above definitions of sufficiency and necessity, we can say what it means to condition a probabilistic causal model \(\langle U, V, F, P(u) \rangle\) on sufficiency and necessity conditions being satisfied.



\begin{definition}[Cause-conditioned distribution]\label{def:cause-conditioned}
    Given a probabilistic causal model \(\langle U, V, F, P(u) \rangle\), a causal set \(\{\vec{C} = \vec{c}\}\), and a set of context nodes \(\vec{T}\) to be kept fixed, a cause-conditioned distribution for the outcome \(\varphi\) is:
    \[P^C_{\vec{C} = \vec{c}, \vec{T}, \varphi}(\vec{u}, \vec{v}) = P_j(\vec{u}, \vec{v}) \cdot \mathbbm{1}\left(N_{\vec{C} = \vec{c}, \vec{T}}(\vec{u}, \varphi)\right) \cdot \mathbbm{1}\left(S_{\vec{C} = \vec{c}, \vec{T}}(\vec{u}, \varphi)\right)\]
\end{definition}


\noindent The intuition here is simple: we take the original distribution $P_j(\vec{u}, \vec{v})$ and we use the corresponding indicators to "reject" cases where one of the conditions fails, effectively setting the resulting probability to zero in such cases.  Note that the cause-conditioned distribution is unnormalized (but can be normalized if need be). 

\subsection{Search over causes and context Nodes}

To compute a cause-conditioned distribution for an outcome, we already need to pass a particular causal set and a set of context nodes. However,  when interested in a causal explanation, one often does not know the exact set of causes or wants to explore various levels of context sensitivity and different context---to achieve this, one needs to carry out a search. This can be achieved by stochastically choosing the causal set and context nodes. To carry out this stochastic choice, we next define a distribution over causal sets and context nodes. Moreover, the probabilities involved in this stochastic search will be "knobs" regulating the extents to which we prefer smaller/larger sets of causes, and to which we want to be context sensitive by keeping parts of the actual setting fixed.

\begin{definition}[Antecedent auxiliary and bias]
    Given a variable \(X\) in the potential causal set \(\{\vec{X} = \vec{x}\}\), let an auxiliary binary variable \(I_X\) denote whether \(X\) is intervened with the value specified in the causal set or not. Given the antecedent bias \(a\), the probability that \(I_X =1\) is \(0.5 - a\) (that is, $a$ is a bias towards preemption of that intervention).
    
\end{definition}

\begin{definition}[Witness auxiliary and bias]
    Given a variable \(W\) in the potential context set \(\vec{W}\), let an auxiliary binary variable \(J_W\) denote whether \(W\) is kept fixed (preempted) at the factual value or not. Given the witness  bias \(w\), the probability that \(J_W=1\)  is \(0.5 + w\) (that is, $w$ is a bias towards preemption of $W$ with the actual value).
\end{definition}

\noindent For an indicator variable $\iota\in \{I, J\}$ and a set of variables \(\vec{B}\), we abbreviate: \(\iota_B = \bigwedge_{A \in \vec{B}} \left(\iota _A = 1 \right)\)    and \(\iota_B^- = \bigwedge_{A\in B} \left(\iota_A = 0\right)\).

\begin{definition}[Interventional setting and its probability distribution] \label{def:interventional_setting} A causal setting is determined by specifying (1) a subset \(C\) of the nodes occurring in a set of potential causes \(\{\vec{X} = \vec{x}\}\) ($C$  will be intervened with the values as specified in that causal set), and (2) a subset \(T\) of the set of potential context nodes \(\vec{W}\) ($T$ will in fact be preempted, i.e.\! intervened to have the factual value). Both can be completely specified by the values of the indicator variables we introduced. Given an antecedent bias \(a\) and witness bias \(w\), the setting probability, i.e. the probability that (exactly) \(C\) is intervened on and (exactly) \(T\) is preempted, $P^{\mathit{set}}_{a,w, \{\vec{X}=\vec {x}\}, \vec{W}}(C, T)$ can be defined in the following terms. Let $A = .5-a$ be the probability that a node is not preempted, i.e., that a proposed intervention is executed. Let $W = 0.5 + w$ be the probability that a witness node is preempted, and also $A' = 1-A$, $W' = 1-W$. In a search, the intention is to intervene on cause candidates stochastically, and in cases where a potential context node is not intervened on as a cause candidate, stochastically preempt it to the observed value (that is, $\vec{X}$ and $\vec{W}$ might overlap in a search, but not any particular $\vec{C}$ and $\vec{T}$). So the corresponding indicators distributions are determined by $I_i \sim Binomial(A), \overset{\sim}{J_i} \sim Binomial(W)$ and $J_i = \overset{\sim}{J_i}(1-I_i)$, which corresponds to the following joint distribution:
\begin{center}
\begin{tabular}{lcc}
\toprule
 & $J_i = 0$ & $J_i = 1$ \\
\midrule
$I_i = 0$ & $A'W'$ & $A'W$ \\
$I_i = 1$ & $A$ & $0$ \\
\bottomrule
\end{tabular}
\end{center}
Then, the probability of any particular interventional setting is:
\begin{align}\label{eq:psetawct}
P^{\mathit{set}}_{a,w, \{\vec{X}=\vec {x}\}, \vec{W}}(C,T) & = P_{a, w}\left(I_c \wedge I^-_{\vec{X} \backslash C}  \wedge  J_T  \wedge J^-_{\vec{W} \backslash T} \right)   % \left( (0.5 - a)^{|C|} (0.5 + a)^{|\vec{X}\backslash C|} \right) \times \left( (0.5 + w)^{|T|}(0.5 - w)^{|\vec{W}\backslash T|} \right)
\end{align}
\noindent 



        
    % \begin{flalign*}
    % P_{a, w}\bigg(\bigwedge_{c \in C} I_c \wedge \bigwedge_{c \in \vec{X} \backslash C} &\neg I_c \bigwedge_{t \in T} J_t \wedge \bigwedge_{t \in \vec{W} \backslash T} \neg J_t\bigg) = \\ & (0.5 + a)^{|C|} (0.5 - a)^{|\vec{X}\backslash C|} \cdot (0.5 + w)^{|T|}(0.5 - a)^{|\vec{W}\backslash T|}
    % \end{flalign*}
\end{definition}

\subsection{Explanations}

Next, we want to be able to weigh cause-conditioned distributions by the interventional setting probabilities. The result is called an explanation-extended distribution.

\begin{definition}[Explanation-extended distribution]
    Given a probabilistic causal model \(\langle \langle U, V, F\rangle, P(u)\rangle\), an outcome \(\varphi\), a set of possible causes \(\{\vec{X} = \vec{x}\}\) and its subset \(C\), a set of context nodes \(\vec{W}\) and its subset \(T\), antecedent bias \(a\) and witness bias \(w\), the explanation-extended distribution can be defined as follows:

    \begin{flalign*}
    P^E_{a, w, \varphi, \{\vec{X} = \vec{x}\}, \vec{W}}(\vec{u}, \vec{v}, C, T) = 
    &P^{\mathit{set}}_{a,w, \{\vec{X}=\vec {x}\}, \vec{W}}(C,T)  P^C_{C, T, \varphi} (\vec{u}, \vec{v})
    \end{flalign*}
\end{definition}


\noindent This means that for any particular interventional setting $(C,T)$, we weight its cause-conditioned distribution by the setting probability distribution corresponding to $(C,T)$.
Note again that the explanation-extended distribution as defined above is unnormalized.

At this point, we have the main components of our approach available. To recap: we start with a joint probability in the baseline model. We multiply it by the necessity and sufficiency indicators to obtain a \emph{cause-conditioned} (unnormalized) distribution. Then we weight it by an interventional setting probability to obtain an \emph{explanation-extended}(unnormalized) distribution. The latter can be used to explicate a range of quantities of interest, to which we will now move. 

\subsection{Probability of causation and related definitions}

The previous subsection defined the notion of an explanation-extended distribution, which can be used to provide explanations for an outcome \(\varphi\). We first define the probability of causation and the degree of responsibility, and elaborate on how these constitute explanations in this setup.

\begin{definition}[Probability of Causation]
    The probability that a causal set \(C\) causes an outcome \(\varphi\) in a probabilistic causal model \(\langle U, V, F, P(u)\rangle\) given the set of potential causes \(\{\vec{X} = \vec{x}\}\), a set of potential context nodes \(\vec{W}\),  and antecedent and witness biases \(a\) and \(w\) is defined as follows:
    \[P(C \rightarrow \varphi) = \sum_{\vec{u}, \vec{v}} \sum_{T \subseteq \vec{W}}P^E_{a, w, \varphi, \{\vec{X} = \vec{x}\}, \vec{W}}(\vec{u}, \vec{v}, C, T) \]
\end{definition}


\noindent That is, to evaluate the "causal power" of a set $\vec{C}$  over an outcome $\varphi$, we (i) intervene with $C$, (ii) marginalize all variables, (iii) marginalize the choice of which context set is kept fixed, and estimate the resulting explanation-extended probability of $\varphi$.


WARNING: THE DISCUSSION OF DEGREE OF RESPONSIBILITY AND EXPLANATIONS THIS WILL BE MODIFIED/EXTEND REVISED, SKIP FOR NOW

\begin{definition}[Degree of Responsibility] The degree of responsibility of a variable setting \(B=b\) in the causal set %\(A = a\)
for an outcome \(\varphi\) is defined as follows:
\[D_R(B \rightarrow \varphi) = \sum_{\vec{u}, \vec{v}} \sum_{T \subseteq \vec{W}} \sum_{C \subseteq \vec{X}, B \in C}P^E_{a, w, \varphi, \{\vec{X} = \vec{x}\}, \vec{W}}(\vec{u}, \vec{v}, C, T)\]
\end{definition}

\begin{definition}[Explanations]
    For an outcome \(\varphi\), a causal explanation is defined as a causal set \(\{\vec{C} = \vec{c}\}\) such that \(P(\vec{C} \rightarrow \varphi) > 0\). The degree of `goodness' of the explanation \(\{\vec{C} = \vec{c}\}\) is defined by its probability of causation.
\end{definition}

\subsection{When variables are continuous}

All the above definitions are compatible with causal models with continuous distributions as well. But these definitions can become infeasible when the outcome \(\varphi\) is a measure-zero event where a continuous random variable exactly takes on a particular value.

This particularly messes with Definition~\ref{def:cause-conditioned} where we create a distribution conditioned on sufficiency and necessity, two measure-zero events. To counter this, we propose that instead of using indicator variables for conditioning, we score them using soft equality and non-equality functions. We first provide the definitions of these score functions and then we provide the extended version of a cause-conditioned distribution.

\begin{definition}[Equality functions]
    An equality function \(f_= : \mathbb{R} \times \mathbb{R} \rightarrow \mathbb{R}\) is a monotonically increasing function of \(d = |r_1 - r_2|\) where \(r_1, r_2\) are its two inputs. 
\end{definition}

\noindent An example of an equality function would be \[f_=(r_1, r_2) = \log{(\mathit{pdf}(|r_1 - r_2|, N(0, \sigma)))}\] where \(\mathit{pdf}(r, N(0, \sigma))\) denotes the pdf value of a zero-centered normal distribution at point \(r\).

\begin{definition}[Inequality functions]
    An inequality function \(f_{\neq} : \mathbb{R} \times \mathbb{R} \rightarrow \mathbb{R}\) is a monotonically decreasing function of \(d =|r_1 - r_2|\) where \(r_1, r_2\) are its two inputs. 
\end{definition}

\noindent An example of an inequality function would be \[f_{\neq}(r_1, r_2) = \log{(2 \cdot \mathit{cdf}(|r_1 - r_2|, N(0, 0.1)) - 1)}\] where \(\mathit{cdf}(r, N(\mu, \sigma))\) denotes the cdf value of a normal distribution at point \(r\).


In our implementation we go a bit further, as we parametrize the functions with a scaling factor $a$ and we want scores to be assigned in the log probability space. So we further require the following:
\begin{itemize}
\item  As $a$ tends to zero, $f_{eq}^a$ tends to the value it takes when $r_1=r_2$ (call it $\tau$, only in that case, and
to negative infinity otherwise, 
\item As $a$ goes to infinity $f_{eq}^a$ tends to $\tau$.

\item As $a$ tends to allowed minimum (which might be finite, depending on  the functional form of $f_{neq}$), the difference between outcomes $f_{neq}^a(r_1,r_1)$ and $f_{neq}^a(r_1,r_2)$ where $r-1 \neq r_2$ goes to negative infinity.

\item As $a$ goes to infinity, $f_{neq}^a(r_1,r_2)$ tends to that of $f_{neq}^a(r_1,r_1)$.

\item For any scale $a$, $f_{neq}^a(r_1,r_2)$ tends to zero
     $d$ tends to infinity, to the minimum of its range as $d$ tends to zero, and doesn't equal the minimum if $d\neq 0$.
\end{itemize}



\begin{definition}[Cause-conditioned distribution (Extended)]
    Given a probabilistic causal model \(\langle U, V, F, P(u) \rangle\), a causal set \(\{\vec{X} = \vec{x}\}\), and a set of context nodes \(\vec{W}\) to be kept fixed, a cause-conditioned distribution for the outcome \(\varphi := (R = r)\) where \(R \in V, r \in \mathit{dom}(R)\) and \(R\) is continuous, if the following hold:
    \[(M, \vec{u}) 
    \vDash [\vec{X} \leftarrow \vec{x}, \vec{W} \leftarrow \vec{w}^\star](R = r_s)\]
    \[(M, \vec{u}) 
    \vDash [\vec{X} \leftarrow \vec{x'}, \vec{W} \leftarrow \vec{w}^\star](R = r_n)\]
    then cause-conditioned distribution can be defined as follows where \(f_=\) and \(f_{\neq}\) are equality and inequality functions respectively:
    \[P^C_{\vec{X} = \vec{x}, \vec{W}, \varphi}(\vec{u}, \vec{v}) = P_j(\vec{u}, \vec{v}) \cdot f_{\neq}(r, r_n) \cdot f_{=}(r, r_s)\]
\end{definition}

\noindent The remaining  definitions remain the same, with only the extended definition of cause-conditioned distribution swapped for the original one.

Taking a step back,  the probabilities of causation, while shaped by the original model itself, are also impacted by what the user does with a few knobs that they might manipulate. Manipulating the probabilities of intervention preemption mirrors user preference for (or against) smaller cause sets. Manipulating the probabilities of witness preemption mirrors how context-sensitive the evaluation is supposed to be. Adding parametrized equality and inequality functions allows the user to decide how steeply the distances between outcomes should be rewarded or penalized by the respective equality/non-equality functions. 


\section{Relation with Actual Causality}

The notion of explanation-extended distribution is not meant to explicate \linebreak Halpern's notion of actual causality; rather, the goal was to combine the intuitions underlying Halpern's notion with those motivating Pearl's notion of probability of necessity and sufficiency to provide a more general tool for causal attribution. Nevertheless, to show that we preserve the spirit of Halpern's definition, in this section, we discuss the connection between actual causality and the notions of explanation-extended distribution. Let's start with the main definition used by Halpern:

\begin{definition} \label{def:ac}
    \(\vec{C} = \vec{c}\) is an actual cause of \(\varphi\) in the causal setting \((M, \vec{u})\) if the following three conditions hold:
    \begin{itemize}
        \item \textbf{Factivity}: \((M, \vec{u}) \vDash (\vec{C} = \vec{c}\,)\) and \((M, \vec{u}) \vDash \varphi\),
        \item \textbf{Context-sensitive Necessity}: There is a set of variables \(\vec{T} \subseteq \vec{V}\) and an alternative setting \(\vec{c'}\) of variables in \(\vec{C}\) such that if \((M, \vec{u}) \vDash \vec{T} = \vec{t^\star}\), then
            \[(M, \vec{u}) \vDash [\vec{C} \leftarrow \vec{c'}, \vec{T} \leftarrow \vec{t^\star}] \neg \varphi,\]
        \item \textbf{Minimality}: \(\vec{C}\) is minimal; there is no strict subset \(\vec{C_{sub}} \subset \vec{C}\) such that \(\vec{C_{sub}} = \vec{c_{sub}}\) satisfies the above two conditions where $\vec{c_{sub}}$ is the restriction of \(\vec{c}\) to the variables in \(\vec{C_{sub}}\). 
    \end{itemize}
\end{definition}


\begin{theorem}\label{th:ac-exp}
    Suppose that $C$ is an actual cause with a witness set $T$ in $(M, \vec{u})$. \linebreak  Take any  $\{\vec{X} = \vec{x}\} \supseteq C$ as the potential cause set, and follow the definition in taking the set of all endogenous variables $V \supseteq T$ as the potential witness set. then for any antecedent bias  $ - 0.5 < a < 0.5$, and  any witness  bias $ - 0.5 < w < 0.5$ we have: \begin{align}
    P^{E}_{a, w , \varphi, \{\vec{X} = \vec{x}\}, \vec{V}}(\vec{u}, \vec{v}, C, T) > 0 \label{eq:pect}\\
    \text{ and therefore } \nonumber\\ 
\sum_{T \subseteq \vec{V}}\, _{a, w , \varphi, \{\vec{X} = \vec{x}\}, \vec{V}}(\vec{u}, \vec{v}, C, T) > 0 \end{align}
\end{theorem}


\begin{proof}  Assuming the antecedent of the theorem, freeze $\vec{u}$, $\vec{c'}$, and $\vec{T}$ as witnesses to the satisfaction of the conditions specified in Definition \ref{def:ac}. For the condition of context-sensitive necessity to be satisfied, $\vec{T} \cap \vec{C}$ need to be disjoint, otherwise the intervention with alternative values $\vec{c'}$ and at the same time actual values $\vec{t^\star}$ would be ill-defined. Further, set $\vec{v}$ to be the values that $\vec{V}$ take in  $(M,u)$, so that $(M, \vec{u}) \models \vec{V} = \vec{v}$.  
We follow the definition of an explanation-extended distribution to decompose \eqref{eq:pect} as follows:
\begin{align}
    P^{\mathit{set}}_{a,w, \{\vec{X}=\vec {x}\}, \vec{V}}(C,T)  P^C_{C, T, \varphi} (\vec{u}, \vec{v}) = \\
      P^{\mathit{set}}_{a,w, \{\vec{X}=\vec {x}\}, \vec{V}}(C,T) \cdot P_j(\vec{u}, \vec{v}) \cdot \mathbbm{1}\left(N_{\vec{C} = \vec{c}, \vec{T}}(\vec{u}, \varphi)\right) \cdot \mathbbm{1}\left(S_{\vec{C} = \vec{c}, \vec{T}}(\vec{u}, \varphi)\right) \label{eq:decomposition}
\end{align}
\noindent The goal now is to show that each of the components is greater than 0.

\begin{itemize}
\item Starting with $P^{\mathit{set}}_{a,w, \{\vec{X}=\vec {x}\}, \vec{W}}(C,T)$, by \eqref{eq:psetawct}, this is identical to \linebreak  $P_{a, w}\left(I_c \wedge I^-_{\vec{X} \backslash C}  \wedge  J_T  \wedge J^-_{\vec{V} \backslash T} \right)$. Given  $\vec{T} \cap \vec{C} = \emptyset$, in light of Definition \ref{def:interventional_setting}, this is the same as 
$\left( (.5-a)^{\mid \vec{C}\mid}   (.5+a)^{\mid \vec{X}\setminus \vec{C}\mid} (.5+w)^{\mid \vec{W}\mid} (.5-w)^{\mid \vec{V}\setminus \vec{W}\mid} )  \right)$, which on the assumption of regularity $>0$.

\item Next, we are looking at $P_j(\vec{u}, \vec{v})$. As $\vect{v}$ are actual values from $(M,\vec{u})$,  $\mathbbm{1}((M, \vec{u}) \vDash \vec{v})=1$ and so by \eqref{eq:joint}  $P_j(\vec{u}, \vec{v}) = P(\vec{u})$. Since $(M, \vec{u})$ are realized by the assumption, $P(\vec{u}) > 0$ and therefore also $P_j(\vec{u}, \vec{v})>0$. 

\item The value of the indicator $\mathbbm{1}\left(N_{\vec{C} = \vec{c}, \vec{T}}(\vec{u}, \varphi)\right)$ is, following \eqref{eq:nec}, determined by there being a $\vec{c'}$ such that 
$(M, \vec{u}) \vDash [\vec{C} \leftarrow \vec{c'}, \vec{T} \leftarrow \vec{t}^\star]\neg \varphi$, which is satisfied by the context-sensitive necessity assumption with the frozen variables as witnesses for the existential claim, so the value of the indicator is 1.

\item For $\mathbbm{1}\left(S_{\vec{C} = \vec{c}, \vec{T}}(\vec{u}, \varphi)\right)$,  note that with exogenous noise $\vec{u}$ held constant, given the first part of the factivity assumption:
\[\left( (M, \vec{u}) 
    \vDash [\vec{C} \leftarrow \vec{c}, \vec{T} \leftarrow \vec{t}^\star]\varphi \right) \Leftrightarrow \left( (M, \vec{u})\models \varphi\right) \] \noindent  Since the right-hand side holds by the second factivity assumption, so does the left-hand side, which yields $\mathbbm{1}\left(S_{\vec{C} = \vec{c}, \vec{T}}(\vec{u}, \varphi)\right) =1$ by ($S_{\vec{C} = \vec{c}, \vec{T}}(\vec{u}, \varphi)$).

\end{itemize}
\end{proof}




\begin{theorem} Moreover, if  
$\sum_{T \subseteq \vec{V}} P^{E}_{a, w, \varphi, C, \vec{V}}(\vec{u}, \vec{v}, C, T) > 0$,  the context-sensitive necessity condition in Definition \ref{def:ac} holds. \label{th:if_pr_then_nec}
\end{theorem}


\begin{proof} Let us fix notation as in the proof of Theorem \ref{th:ac-exp}.  Since the probabilities are non-negative, for the assumption to hold, there has to be at least one $T$ for which \eqref{eq:pect} holds, and so we are now looking again at the decomposition in \eqref{eq:decomposition}. However, now, by the assumption, the whole composite is positive, and therefore so is $\mathbbm{1}\left(N_{\vec{C} = \vec{c}, \vec{T}}(\vec{u}, \varphi)\right)$, which by \eqref{eq:nec} means that there are $\vec{c'}$ and $\vec{T}$ such that: 
\[
    (M, \vec{u}) \vDash [\vec{C} \leftarrow \vec{c'}, \vec{T} \leftarrow \vec{t}^\star]\neg \varphi.\] These are the witnesses needed to validate the existential claim involved in Definition \eqref{def:ac}.
\end{proof}


\pagebreak 


With minimality, the situation is somewhat more complicated. For one thing, our choice of antecedent and witness preemption biases will have an impact on probability-preferred set sizes. On the other hand, the way that \linebreak $P^{\mathit{set}}_{a,w, \{\vec{X}=\vec {x}\}, \vec{V}}(C,T)$ is defined makes sets comparable in light of their cardinalities, whereas the minimality Definition \ref{def:ac}  is formulated in terms of the subset relation, which is bound to produce some misalignment. Some useful observations can be made, though. 


\begin{theorem} \label{th:local_max_to_ac} Write $\Phi (C, T)$ for $P^{E}_{a, w, \varphi, C, \vec{V}}(\vec{u}, \vec{v}, C, T)$. Consider the set of optimizers  $\Gamma = \arg\max_{\vec{C} \subseteq \{\vec{X} = \vec{x}\}, \vec{T} \subseteq \vec{V}} \Phi(C, T)$.  Assume $(C^\dagger, T^\dagger)\in \Gamma$ with $\Phi(C^\dagger, T^\dagger)>0$ and assume further that factivity is satisfied for $C^\dagger$ (as in Definition \ref{def:ac}). Let the witness bias be $w=0$ and let the antecedent preemption bias be $a>0$ (that is, we positively prefer antecedent preemption, thus building preference of smaller cause sets into our search). Then $C^\dagger$ is an actual cause of $\varphi$ according to Definition \ref{def:ac}.
\end{theorem}

\begin{proof}
We already assumed factivity, and by Theorem \ref{th:if_pr_then_nec} we know the context-sensitive necessity condition is satisfied. So the only reason why $C^\dagger$ could fail to be an actual cause is minimality. We will assume minimality indeed fails and reason to the conclusion that in that case $(C^\dagger, T^\dagger)$ couldn't belong to $\Gamma$ in the first place.


So suppose there is a $C^\sharp \subset C^\dagger$ which is an actual cause with a witness $T^\sharp$. $\Phi(C^\dagger, T^\dagger)$ and $\Phi(C^\sharp, T^\sharp)$ are, respectively, as follows:
\begin{align}
P^{\mathit{set}}_{a,w, \{\vec{C^\dagger}=\vec {c^\dagger}\}, \vec{V}}(C^\dagger,T^\dagger) \cdot P_j(\vec{u}, \vec{v}) \cdot \mathbbm{1}\left(N_{\vec{C^\dagger} = \vec{c^\dagger}, \vec{T^\dagger}}(\vec{u}, \varphi)\right) \cdot \mathbbm{1}\left(S_{\vec{C^\dagger} = \vec{c^\dagger}, \vec{T^\dagger}}(\vec{u}, \varphi)\right)  \label{eq:comparisondagger}\\
P^{\mathit{set}}_{a,w, \{\vec{C^\sharp}=\vec {c^\sharp}\}, \vec{V}}(C^\sharp,T^\sharp) \cdot P_j(\vec{u}, \vec{v}) \cdot \mathbbm{1}\left(N_{\vec{C^\sharp} = \vec{c^\sharp}, \vec{T^\sharp}}(\vec{u}, \varphi)\right) \cdot \mathbbm{1}\left(S_{\vec{C^\sharp} = \vec{c^\sharp}, \vec{T^\sharp}}(\vec{u}, \varphi)\right) \label{eq:comparisonsharp}
\end{align}

\noindent In comparing them, we can drop  $P_j(\vec{u}, \vec{v})$ as identical. Moreover, indicators in \eqref{eq:comparisondagger} are 1s on the assumption that $\Phi(C^\dagger, T^\dagger)>0$ and indicators in \eqref{eq:comparisonsharp} are 1s by Theorem \ref{th:ac-exp} and the assumption that $C^\sharp$ is an actual cause with witness $T^\sharp$.  So all that is left to compare is $P^{\mathit{set}}_{a,w, \{\vec{X}=\vec {x}\}, \vec{V}}(C^\dagger,T^\dagger)$ and  $P^{\mathit{set}}_{a,w, \{\vec{X}=\vec {x}\}, \vec{V}}(C^\sharp,T^\sharp)$.

\begin{align}
P^{\mathit{set}}_{a,w, \{\vec{X}  =\vec {x}\}, \vec{V}}(C^\dagger,T^\dagger) & =  P_{a, w}\left(I_{C^\dagger} \wedge I^-_{\vec{X} \backslash C^\dagger}\right) \cdot \underbrace{P_{a,w} \left( J_{T^\dagger}  \wedge J^-_{\vec{V} \backslash T^\dagger} \right)}_{ = .5^{\vert \vec{V} \vert} \text{ by $w=0$}}\\
P^{\mathit{set}}_{a,w, \{\vec{X}=\vec {x}\}, \vec{V}}(C^\sharp,T^\sharp) & =  P_{a, w}\left(I_{C^\sharp} \wedge I^-_{\vec{X} \backslash C^\sharp}  \right) \cdot \overbrace{P_{a,w}  \left(  J_{T^\sharp}  \wedge J^-_{\vec{V} \backslash T^\sharp} \right)}
\end{align}
\noindent We are now left comparing $(.5-a)^{\vert C^\dagger\vert } (.5+a)^{\vert \vec{X}\setminus{C^\dagger}\vert}$ with $(.5-a)^{\vert C^\sharp\vert } (.5+a)^{\vert \vec{X}\setminus{C^\sharp}\vert}$. But by the assumption that $C^\sharp \subset C^\dagger$, we know $\vert C^\dagger\vert > \vert C^\sharp\vert$ and $\vert \vec{X}\setminus{C^\dagger}\vert < \vert \vec{X}\setminus{C^\sharp}\vert$. Thus $P^{\mathit{set}}_{a,w, \{\vec{X}=\vec {x}\}, \vec{V}}(C^\sharp,T^\sharp) > P^{\mathit{set}}_{a,w, \{\vec{X}  =\vec {x}\}, \vec{V}}(C^\dagger,T^\dagger)$ and so $\Phi(C^\sharp, T^\sharp) > \Phi(C^\dagger, T^\dagger)$, which  contradicts the assumption that $(C^\dagger, T^\dagger) \in \Gamma$.
\end{proof}




\begin{corollary} Let's keep the notation from Theorem \ref{th:local_max_to_ac}. Consider the set   of optimizers  $\Delta = \arg\max_{\vec{C} \subseteq \{\vec{X} = \vec{x}\}} \sum_{\vec{T} \subseteq \vec{V}} \Phi(C, T)$.  Assume $C^\dagger\in \Delta$ with \[\sum_{\vec{T} \subseteq \vec{V}} \Phi(C^\dagger, T)>0\] and assume still that factivity is satisfied for $C^\dagger$. Keep $w=0$ and $a>0$  as in Theorem \ref{th:local_max_to_ac}. Then $C^\dagger$ is an actual cause of $\varphi$ according to Definition \ref{def:ac}.
\end{corollary}


\begin{proof}
The assumptions entail there is a $T^\dagger\in \mathcal{P}(\vec{V})$ such that $(C^\dagger, T^\dagger) \in \Gamma$ (as defined before) and $\Phi(C^\dagger, T^\dagger)>0$, because $\{(C^\dagger, T)\vert T\subseteq \vec{V}\} \supseteq \Gamma$. This gives a witness set that together with $C^\dagger$ satisfies the conditions of Theorem \ref{th:local_max_to_ac}, so by that theorem it follows that $C^\dagger$ is an actual cause.
\end{proof}

\newpage

Now we can return to the misalignment we have mentioned and elaborate on the nuance. One can imagine a reasoning mode when one looks for causes of an outcome by inspecting probabilities of causation $P(C \rightarrow \varphi)$ not knowing what the values of $\vec{u}$ are, but knowing that $\varphi$ occurred. It is worth pointing out that  $\arg\max_{\vec{C} \subseteq \vec{X}}P(\vec{C} \rightarrow \varphi)$ might fail to include some $\vec{C}$ such that for some setting $\vec{u}$ $\vec{C}$ is nevertheless an actual cause of $\varphi$.


\begin{observation} Define $\Xi = \arg\max_{\vec{C} \subseteq {X}} P(C \rightarrow \varphi)$. It is not the case that if $C^\dagger$ is an actual cause in $(M, \vec{u})$ for some $\vec{u}$, $C^\dagger \in \Xi$.
\end{observation}

\begin{proof}
Consider a very simple model with two parent binary nodes $A, B$ and a binary child node $Y$, with uniform distribution over the states of $A, B$. Let $C  =  \mathbbm{1}\left( (M, \vec{u}) \models A \vee B\right)$. Consider $C^\dagger = \{A = 1, B =1\}$ and $C^{\sharp} = \{A=1\}$. Let $w$ be 0, so that all witness sets are equally probable. Let $u_{ij}$ represent a setting where $A=i, B= j$.
\begin{itemize}
\item In $u_{11}$, $C^\dagger$ is an actual cause of $\varphi := Y=1$. Setting the values of $A, B$ to the only alternative results in $\neg\varphi$, no change in any of the subsets to $C^\dagger$ does so, and activity is satisfied. 
\item There is effectively only one non-zero-measure witness set for $C^\dagger$, which is $\emptyset$. With that witness set, necessity indicators and the sufficiency indicators are always 1 (always, setting $A$ and $b$ to ones results in $Y=1$, and always, setting them to zeros results in $Y=0$. The probability of each scenario is $1/4$ (there are four such settings), and the probability of that causal setting is $(.5-a)^2$. So, $P(C^\dagger \rightarrow \varphi) = 4 \times ((.5-a)^2 \times 1/4)$.


\item in $u_{10}$, $C^\sharp$ is an actual cause of $\varphi$. There are two witness sets, $\emptyset$ and $\{B\}$, each equally likely. The sufficiency indicator for $C^\sharp$ is 1 everywhere, but the necessity indicator for $C^\sharp$ is 1 in $u_{10}$ and $u_{00}$,  so we only sum over those. So we multiply by 2 because of 2 witness sets, and by 2 again, because there are two settings of $\vec{u}$ where both indicators are 1.
This means we multiply by 4 as before. However, this time the probability of the causal setting is $(.5-a)(.5+a)$, so together we have that $P(C^\sharp \rightarrow \varphi) = 4( (.5-a)(.5+a) \times 1/4) $
\end{itemize}
This means that the difference between these two probabilities of causation will boil down to differences resulting from the choice of $a$. On rare occasions, they will be identical, but for many choices they will not, and only one of $C^\dagger$, $C^\sharp$ will end up in $\Xi$ (see Figure \ref{fig:simple_max}).

\begin{figure}
    \centering
    \includegraphics[width=0.8\linewidth]{figures/set_probabilities.png}
    \caption{Probabilities of causation in a simple three-node model; for each set, there exists a setting $\vec{u}$ in which that set is an actual cause of $Y=1$, but whether the probability of causation $P(Set \rightarrow Y=1)$ is maximized depends on the choice of the antecedent preemption bias. This ambiguity does disappear once we identify which $\vec{u}$ are actual. But this also suggests that inspecting only the maximizers of the probability of causation is not a robust strategy of identifying actual causes across a range of possible values of the exogenous nodes.}
    \label{fig:simple_max}
\end{figure}


\end{proof}


\pagebreak 

\begin{observation} Consider 
\begin{align*}\Xi & = \left\{C \vert \left(\sum_{\vec{T} \subseteq \vec{V}} \Phi(C, T)>0\right),  \neg \exists C' \subset C \left( \sum_{\vec{T} \subseteq \vec{V}} \Phi(C', T)\right) \right\}
\end{align*}

\begin{enumerate} 
\item $\Delta \subseteq \Xi$.
\item If $C^\dagger \in \Xi$, $C^\dagger$ is an actual cause. 
\end{enumerate}
\end{observation}



\begin{proof}
TODO
\end{proof}


% \subsection{Factivity}

% Our definition of explanation-extended distribution does not enforce the criteria of factivity but that can be easily achieved by conditioning. We also believe that the factivity is a restricting condition since we do not wish to limit our causal reasoning to a particular factual world. We wish to generalize to other possible worlds as well.



% \section{Questions}
% \begin{itemize}
%     \item These definitions do not exactly coincide with Halpern and Pearl's definitions. How do we want to go about justifying these definitions?
%     \item Are there any other quantities that people think we can define above? 
% \end{itemize}














\section{Discussion}

\subsection{Relation to explainable AI and related desiderata}




%TODO: causal explanation of reality vs. causal explanation of an algorithmic decision

\cite{verma2022recourses} in a wide survey put forward a list of potentially desirable features for counterfactual explanations. Let's go over them and discuss how our approach fares.

\begin{description}
    \item[Actionability] A recommended counterfactual should not "change immutable features". \cite{verma2022recourses} suggest factoring in how difficult it is to change the value of a given feature into the loss function when generating a counterfactual explanation. As a research challenge, they insist that counterfactual explanations should capture the applicant’s preferences.  Similarly, \cite{Karimi2021} recommend  shifting the focus from explanations to interventions.

    While we agree that the information provided by an explanation should be useful for guiding action and that not all counterfactual explanations are, we don't think actionability and causal responsibility should be conflated. It is one thing to understand the system and the problem at hand, it is another to decide how to go about handling it. Suppose one wants to use counterfactual explanations to help lease applicants, such as Alice or Bob. If the method already has some notion of ease of change baked into the cause identification procedure and never recommends difficult changes, it will never identify gender as a causal factor, as changing one's gender to obtain a lease is a radical move. But this will prevent users from noticing systemic injustice or structural problems. Instead of helping the users notice them, an algorithm that conflates actionability with responsibility would just assist the users in gaming the system as is, without allowing them to reflect on the role features that are hard to change play in it. Moreover, typical phenomena used as a proxy for actionability in optimization, such as counting the number of features that would need to be changed, or some distance measure telling how much a feature would need to change to change the outcome, are poor replacements for realistic actionability considerations. 

   That said, there’s nothing stopping us from combining a probabilistic identification of factors responsible for a given outcome with a decision-theoretic layer. This layer could explicitly account for the costs of modifying various features and optimize for the least expensive actions needed to change the outcome if that is what one's after. In principle, users could integrate these layers to develop more efficient optimization algorithms when actionable causes are their focus. However, this should be done with full awareness of the underlying assumptions and the potential mismatch between factors that are responsible and those that are actionable.


\item[Model agnosticism] A model agnostic method is more widely applicable. Our approach is meant to be used with ChiRho causal models. However, it can be applied to a model-agnostic predictor function based on a black-box model. However, the price of doing so will be the resulting blindness to the causal relationship between the features involved. Our implementation, however, allows one to specify a causal model for the features involved and compose it with the predictor function, regaining some sensitivity to the underlying causal structure. 

\item[Multiple counterfactuals] Many XAI algorithms mostly return a single counterfactual for a
given input data point. Our implementation  allows for the estimation of multiple counterfactuals in one run. 



\item[Sparsity] A method should prefer simpler explanations. Our implementation allows for basing causal search towards smaller sets of features.

\item[Causal constraint satisfaction] The method should be realistic in terms of causal dependencies between features. As soon as our method is applied to a model that contains a representation of the causal structure for the relevant features, our search for explanation satisfies causal constraints.


\end{description}








\bibliographystyle{plain}
\bibliography{references}

\end{document}
